\chapter{Tab 2 · SOP: the laboratory's rules as a profile}\label{ch:sop}

The separate Tab~3 (Results) uses the same active profile. It is opened automatically after files are read, so the
interpretation table follows the loading step while the editable profile remains in this tab.

\section{The SOP profile}
Interpretation rules differ between laboratories and methods. The page keeps them in one \textbf{SOP profile}, a JSON document
(schema \texttt{qpcr-sop-profile/1}) that can be edited on the page, saved, locked and reopened. The stored Cq values and calls
are never changed; every outcome says which rule produced it, and every export carries the profile's name, version and SHA-256.

\begin{center}\small
\begin{tabular}{@{}lp{11.6cm}@{}}\toprule
Section & Content\\\midrule
1 · Identity & name, version, author or approver, assay, effective date, notes\\
2 · Targets & target patterns (exact, wildcard or regular expression), kind (target, reference, internal control), Cq cut-off, late-signal limit, minimum quantity, outcome in the late zone\\
3 · Controls & how controls are recognised (by role or by sample name), expected result (negative, positive with a Cq range, standard), purpose, material, minimum replicate wells per group, minimum groups per run, action on failure\\
4 · Replicates \& IC & minimum replicates, replicates needed for a positive, maximum SD, internal-control shift and fixed reference Cq\\
5 · Run acceptance & integrity, missing or failed controls, standard-curve R$^2$ and efficiency, NTC gap, edit delay, calibration, zone spread, run state\\
6 · Outcomes & labels, colours and the report sentence of each outcome\\
7 · Graph settings & which SOP lines are drawn\\
8 · Control charts & per-chart settings of the Control panel (\texttt{controlCharts})\\\bottomrule
\end{tabular}
\end{center}

The SC figures use the generated review profile \emph{SC screening -- SOP-06 laboratory default}, version 1.0. It is review-only until the laboratory approves the mapping and the unresolved comparison direction in SOP-06 section 6.2(c5). The profile recognises 35S and T-NOS targets, HMG or LEC reference targets, 415B and 410CP LOD controls, NTC, process blanks and negative maize or soy controls. This abbreviated JSON shows the profile shape; values are synthetic configuration examples:

\begin{lstlisting}[style=vstyle,language=json, numbers=left,firstnumber=1]
{ "schema": "qpcr-sop-profile/1", "name": "SC screening -- SOP-06 laboratory default", "version": "1.0",
  "author": "Demo QA manager",
  "targets": [ {"match": "HMG|hmg|lec|Le1|ADH|adh1|SSIIb|zSSIIb|Cru|cruA|GLP", "kind": "reference",
                "cqMax": 40, "cqLate": 36, "lateOutcome": "Inconclusive"},
               {"match": "*", "kind": "target", "cqMax": 40, "cqLate": 38, "lateOutcome": "Repeat"} ],
  "controls": [ {"name": "415B LOD control", "matchBy": "name", "match": "/^415B(?=\\b|\\d)/i",
                 "targets": "35S|TNOS|T-NOS", "expect": "positive", "purpose": "LOD control",
                 "cqLo": 29, "cqHi": 36, "minReplicates": 2, "minPerRun": 0, "onFail": "review"},
                {"name": "Extraction blank", "matchBy": "name", "match": "/^(BLANK|BLK)\\s*EX/i",
                 "expect": "negative", "purpose": "extraction blank", "minReplicates": 1,
                 "minPerRun": 0, "onFail": "review"} ],
  "replicates": {"min": 2, "positiveMin": 2, "partialOutcome": "Repeat", "maxSd": 0.5, "sdOutcome": "Repeat"},
  "run": {"integrity": "reject", "controlsFail": "reject", "editDelayHours": 72, "editAction": "review", "…": "…"},
  "controlCharts": {"cool_rate": {"x": "date", "baseFrom": "2025-05-12"}, "transition": {"x": "date"}},
  "sha256": "…" }
\end{lstlisting}

\par\medskip\noindent{\small\textbf{Listing}~\texttt{sopBase()}\hfill\textcolor{gray}{HTML lines 4095--4111}}\par\nopagebreak
\lstinputlisting[style=vstyle,language=JavaScript,firstnumber=4095]{code/sopBase.js}


\par\medskip\noindent{\small\textbf{Listing}~\texttt{SOP\_PRESETS()}\hfill\textcolor{gray}{HTML lines 4112--4139}}\par\nopagebreak
\lstinputlisting[style=vstyle,language=JavaScript,firstnumber=4112]{code/SOP_PRESETS.js}



\textbf{Which starting point the page offers.} \texttt{Start from} now lists only \emph{SC screening — SOP-06 laboratory
default} (\texttt{SOP\_PRESETS.sc}), and a session with no stored profile starts from it. The earlier generic, GMO and
forensic templates remain in the source as scaffolds used by the review mapping (\texttt{sopLabScreening}), but they are not
offered in the selector and a stored profile recognised as one of those templates is replaced by the SC default when only
\texttt{SC\_} experiments are loaded. Opening a file with the review schema
\texttt{qpcr-control-mapping-proposal/v1} converts it to the SC profile shape; the raw names, stored Cq values, calls and
dates are untouched.

\section{Integrity of the profile}
The SHA-256 is computed over a canonical form of the profile (keys sorted, the fields \texttt{sha256} and \texttt{locked} left
out). Any change of any rule, including a control-chart setting, changes the hash, which is printed on every graph and export.

\par\medskip\noindent{\small\textbf{Listing}~\texttt{sopCanonical()}\hfill\textcolor{gray}{HTML lines 4239--4242}}\par\nopagebreak
\lstinputlisting[style=vstyle,language=JavaScript,firstnumber=4239]{code/sopCanonical.js}


\par\medskip\noindent{\small\textbf{Listing}~\texttt{sopHash()}\hfill\textcolor{gray}{HTML lines 4243--4243}}\par\nopagebreak
\lstinputlisting[style=vstyle,language=JavaScript,firstnumber=4243]{code/sopHash.js}



\figui[0.8]{ui_sop.png}{SOP \& results tab: profile header with version and SHA-256, sections 1--2 of the editor; sections 3--8 are folded.}{ui_sop}

\section{Evaluating a run}
\texttt{sopEvaluateRun()} groups the wells by sample and target, recognises controls, applies cut-offs, late-signal limits,
replicate rules and the internal control, and then the run-acceptance criteria. A run that fails a criterion set to
\emph{reject} turns all its sample results into \emph{Invalid run}. In the synthetic data set the planted contaminated extraction
blank of run 022 does exactly this (Figure~\ref{fig:g_LCX_curves}).

\par\medskip\noindent{\small\textbf{Listing}~\texttt{sopMatcher()}\hfill\textcolor{gray}{HTML lines 4248--4254}}\par\nopagebreak
\lstinputlisting[style=vstyle,language=JavaScript,firstnumber=4248]{code/sopMatcher.js}


\par\medskip\noindent{\small\textbf{Listing}~\texttt{sopEvaluateRun() — opening part}\hfill\textcolor{gray}{HTML lines 4434--4570}}\par\nopagebreak
\lstinputlisting[style=vstyle,language=JavaScript,firstnumber=4434]{code/sopEvaluateRun.js}



\section{Dates and the run timeline}
Every dated event in a file (experiment created, run started and ended, analyses created and modified, last change; for
.eds instrument also the ZIP entry dates, instrument log and calibrations) becomes a timeline. The edit delay and the run duration
used by the SOP and the control panel come from it.

\par\medskip\noindent{\small\textbf{Listing}~\texttt{sopTime()}\hfill\textcolor{gray}{HTML lines 4370--4374}}\par\nopagebreak
\lstinputlisting[style=vstyle,language=JavaScript,firstnumber=4370]{code/sopTime.js}


\par\medskip\noindent{\small\textbf{Listing}~\texttt{runTimeline()}\hfill\textcolor{gray}{HTML lines 4387--4407}}\par\nopagebreak
\lstinputlisting[style=vstyle,language=JavaScript,firstnumber=4387]{code/runTimeline.js}


\par\medskip\noindent{\small\textbf{Listing}~\texttt{runTimes()}\hfill\textcolor{gray}{HTML lines 4408--4415}}\par\nopagebreak
\lstinputlisting[style=vstyle,language=JavaScript,firstnumber=4408]{code/runTimes.js}



\section{SOP findings in the forensic log}
\par\medskip\noindent{\small\textbf{Listing}~\texttt{sopEvents()}\hfill\textcolor{gray}{HTML lines 4623--4635}}\par\nopagebreak
\lstinputlisting[style=vstyle,language=JavaScript,firstnumber=4623]{code/sopEvents.js}




